\section{Protokół wyrównania budżetu obliczeniowego} \label{sec:budget}

Rzetelne porównanie wariantów algorytmów wymaga przydzielenia każdemu z nich identycznego budżetu obliczeniowego. Wyrównanie wyłącznie według liczby wywołań funkcji fitness jest niewystarczające --- w schemacie koewolucyjnym jedno wywołanie ewaluatora ocenia pulę $2\mu$ osobników metodą każdy z każdym, generując $\binom{2\mu}{2} \times 9$ wywołań symulatora (9 kombinacji talii), natomiast w SHADE z Hall of Fame analogiczny koszt wynosi $N_{\text{opp}} \times 9$. Wyrównanie tylko po liczbie wywołań funkcji przystosowania prowadziłoby do przydzielenia bardzo różnej mocy obliczeniowej poszczególnym konfiguracjom.

\subsection{Jednostka budżetowa}

Podstawową jednostką obliczeniową jest jedno wywołanie symulatora SabberStone, czyli pojedyncze uruchomienie symulatora rozgrywające \texttt{NUM\_GAMES = 20} partii między dwoma wektorami wag w ramach jednej kombinacji talii. Liczba wywołań symulatora zależy wyłącznie od struktury algorytmu --- rozmiaru populacji, schematu ewaluacji i liczby przeciwników --- i jest niezależna od wartości \texttt{NUM\_GAMES}.

\subsection{Liczba przeciwników w SHADE z Hall of Fame}

W schemacie koewolucyjnym każda generacja łączy $\mu$ rodziców i $\mu$ potomków w pulę $2\mu$ osobników, a każdy osobnik mierzy się z dokładnie $2\mu - 1$ przeciwnikami. Aby zapewnić porównywalną różnorodność przeciwników w SHADE z Hall of Fame, liczbę stałych przeciwników ustalono na:
\begin{equation}
    N_{\text{opp}} = 2\mu - 1
    \label{eq:nopp}
\end{equation}
Wartości dla testowanych rozmiarów populacji zestawiono w tabeli~\ref{tab:nopp}.

\begin{table}[h]
\centering
\caption{Liczba stałych przeciwników $N_{\text{opp}}$ w zależności od rozmiaru populacji.}
\label{tab:nopp}
\begin{tabular}{cc}
\toprule
$\mu$ & $N_{\text{opp}} = 2\mu - 1$ \\
\midrule
5  & 9  \\
10 & 19 \\
15 & 29 \\
\bottomrule
\end{tabular}
\end{table}

\subsection{Wzory na liczbę wywołań symulatora}

\paragraph{Schemat koewolucyjny.}

W generacji~0 pula zawiera wyłącznie $\mu$ osobników początkowych --- populacja potomna jeszcze nie istnieje w pierwszym kroku. W kolejnych generacjach pula liczy $2\mu$ osobników. Koszt generacji~0 oraz koszt jednej generacji ewolucyjnej wynoszą odpowiednio:
\begin{equation}
    B_0 = \frac{\mu(\mu-1)}{2} \times 9, \qquad
    B_{\text{evo}} = \mu(2\mu - 1) \times 9
    \label{eq:bevo}
\end{equation}
Całkowita liczba wywołań symulatora:
\begin{equation}
    B_{\text{koew}} = B_0 + G_{\text{evo}} \times B_{\text{evo}},
    \qquad G_{\text{evo}} = \frac{\texttt{MAX\_EVALUATIONS}}{\mu} - 1
    \label{eq:btotal_ga}
\end{equation}

\paragraph{SHADE z Hall of Fame.}

Każdy osobnik oceniany jest względem $N_{\text{opp}}$ stałych przeciwników we wszystkich kombinacjach talii. Co $K_{\text{HOF}} = 5$ generacji mechanizm Hall of Fame ponownie ewaluuje $\mu$ rodziców, a koszt tej reewaluacji nie jest wliczany do licznika \texttt{check\_evals}. Łączna liczba wywołań funkcji fitness wynosi:
\begin{equation}
    F_{\text{total}} = \mu + \texttt{check\_evals}
    + \left\lfloor \frac{\texttt{check\_evals}}
      {I_{\text{HOF}} \cdot \mu} \right\rfloor \times \mu
    \label{eq:ftotal}
\end{equation}
Łączna liczba wywołań symulatora:
\begin{equation}
    B_{\text{HoF}} = F_{\text{total}} \times N_{\text{opp}} \times 9
    \label{eq:btotal_pure}
\end{equation}

\subsection{Budżet bazowy}

Punkt odniesienia stanowi eksperyment García-Sáncheza~\cite{sanchez2020} przy $\mu = 10$ i \texttt{MAX\_EVALUATIONS = 1000}:
\begin{equation}
    B^* = \frac{10 \cdot 9}{2} \times 9
        + 99 \times 10 \times 19 \times 9
    = 405 + 169 290
    = \mathbf{169 695 \text{ wywołań symulatora}}
    \label{eq:bstar}
\end{equation}
co odpowiada $3 393 900$ rozegranych partii Hearthstone ($B^* \times 20$).

\subsection{Wyznaczanie wartości \texttt{check\_evals}}

Dla SHADE koewolucyjnego wartość \texttt{check\_evals} wyznacza się analitycznie z równania~\eqref{eq:btotal_ga}:
\begin{equation}
    \texttt{check\_evals}(\mu) =
    \frac{\dfrac{B^*}{9} - \dfrac{\mu(\mu-1)}{2}}{2\mu - 1}
    \label{eq:check_coev}
\end{equation}
Dla SHADE z Hall of Fame wartość wyznaczana jest iteracyjnie jako największa liczba całkowita spełniająca warunek $F_{\text{total}} \leq F^*_{\text{total}}$, gdzie $F^*_{\text{total}} = B^* / (9 \times N_{\text{opp}})$.

\subsection{Wyniki wyrównania}

Tabele~\ref{tab:budget_coev} i~\ref{tab:budget_hof} zestawiają wyznaczone wartości \texttt{check\_evals} wraz z uzyskanymi budżetami i odchyleniami od $B^*$. Wszystkie konfiguracje mieszczą się w granicach $0{,}80\%$ budżetu docelowego, a odchylenia są zawsze ujemne lub zerowe --- budżet nigdy nie jest przekraczany. Wynikają one ze skokowego charakteru kosztu generacji --- przy dużym $\mu$ żadna całkowita liczba generacji nie daje dokładnie wartości $B^*$.

\begin{table}[!h]
\centering
\caption{Wyrównanie budżetu dla SHADE koewolucyjnego.}
\label{tab:budget_coev}
\begin{tabular}{cccccc}
\toprule
$\mu$ & \texttt{check\_evals} & $G_{\text{evo}}$ & $B_0$ &
$B_{\text{evo}}$/gen & Łącznie $B$ \\
\midrule
5  & 2090 & 418 & 90    & 405     & 169 380 \\
10 & 990  & 99  & 405   & 1 710 & 169 695 \\
15 & 645  & 43  & 945   & 3 915 & 169 290 \\
\bottomrule
\end{tabular}
\end{table}

\begin{table}[h]
\centering
\caption{Wyrównanie budżetu dla SHADE z Hall of Fame.}
\label{tab:budget_hof}
\begin{tabular}{ccccccc}
\toprule
$\mu$ & $N_{\text{opp}}$ & \texttt{check\_evals} &
Aktualizacje HoF & $F_{\text{total}}$ & Łącznie $B$ &
$\Delta$ od $B^*$ \\
\midrule
5  & 9  & 1745 & 69 & 2 095 & 169 695 & $0$        \\
10 & 19 & 820  & 16 & 990     & 169 290 & $-405$     \\
15 & 29 & 525  & 7  & 645     & 168 345 & $-1 350$ \\
\bottomrule
\end{tabular}
\end{table}

Odchylenia są nieistotne z punktu widzenia analizy porównawczej --- margines $0{,}80\%$ jest pomijalny wobec wariancji wynikającej z losowości przebiegu partii Hearthstone przy skończonej próbie \texttt{NUM\_GAMES = 20}.
